\section{Evaluation and Error Analysis}
\label{sec:eval}
We first evaluate the archived web-client path and then use explicit gold fields as a controlled diagnostic.
A separate regression analysis concerns the revised policy.
The systems, scorers, and input conditions are kept distinct because changing a scorer can alter a count without changing any sentence.

\subsection{Archived protocol and scoring}
The archived set has 184 monolingual scenarios: 46 selections in each language, comprising ten object requests, ten refusals, eight feelings, eight negated feelings, six action--object selections, and four help or call selections.
Nine selections duplicate few-shot examples and are excluded from the primary analysis, leaving 175 scenarios: 46 English, 45 Hindi, 42 Gujarati, and 42 Marathi.
Eight Hinglish scenarios are exploratory and are not pooled with these counts.
The archived run used one NVIDIA GTX~1650 with 4\,GB memory.

The lexical scorer requires every selected concept, correct gold polarity and act, no detected unselected concept, and an on-task sentence in the target script.
Concept retention has a different denominator: 215 required content-concept instances, rather than 175 scenarios.
Wrong act, missing act, added content, and off-task flags can overlap.
An off-task refusal is a model declining to answer, such as “I cannot help with that”; a legitimate user refusal, such as “I do not want pizza,” is on-task.
Lexical success does not establish grammaticality, meaning in context, or usefulness to AAC users.

The comparison includes literal labels, one unconstrained prompt to the same model, the archived cascade candidate before checking, the deployed policy, and the matched template.
The candidate and final text come from the same recorded response, isolating the joint effect of the gate, retry, and fallback.
The unconstrained prompt requests one sentence without a ledger, examples, refinement, or verification, using greedy decoding with a 75-token output cap; its 8/175 score diagnoses that prompt, rather than all one-pass methods.
Paired pooled intervals resample 46 underlying selection groups across their language versions, using 10,000 bootstrap resamples.

\begin{table*}[t]
\centering\small
\setlength{\tabcolsep}{5pt}
\begin{tabular}{@{}L{5.8cm}ccccc@{}}
\toprule
System / condition & English & Hindi & Gujarati & Marathi & Pooled \\
\midrule
\multicolumn{6}{@{}l}{\textit{A. Archived web client; scorer v1}}\\
Unconstrained Qwen (diagnostic) & 2/46 & 1/45 & 4/42 & 1/42 & 8/175 \\
Cascade before checking & 3/46 & 24/45 & 20/42 & 18/42 & 65/175 \\
Gate + retry + fallback & 22/46 & 31/45 & 21/42 & 23/42 & 97/175 \\
Matched template baseline & 33/46 & 31/45 & 27/42 & 27/42 & 118/175 \\
\midrule
\multicolumn{6}{@{}l}{\textit{B. Archived explicit gold act and polarity; scorer v1}}\\
Cascade before checking & 5/46 & 26/45 & 16/42 & 24/42 & 71/175 \\
Gate + retry + fallback & 33/46 & 41/45 & 26/42 & 32/42 & 132/175 \\
Matched template baseline & 46/46 & 45/45 & 34/42 & 42/42 & 167/175 \\
\midrule
\multicolumn{6}{@{}l}{\textit{C. Revised browser path; scorer v2; regression only}}\\
Revised cascade before checking & 6/46 & 29/45 & 27/42 & 26/42 & 88/175 \\
Revised gate + retry + checked fallback & 46/46 & 45/45 & 42/42 & 42/42 & 175/175 \\
Repaired template baseline & 46/46 & 45/45 & 42/42 & 42/42 & 175/175 \\
\bottomrule
\end{tabular}
\caption{Scenario success as counts over the same 175 selections. Panels A and B reproduce the archived lexical scorer; panel C is a scorer-v2 regression comparison on scenarios that informed the repairs. Scores across scorers are not directly comparable. English is a language reference; Hindi, Gujarati, and Marathi evaluate language extensions. The template is the method baseline in every language.}
\label{tab:results}
\end{table*}


\subsection{Web-client results and gold-field diagnostic}
With actual browser-inferred fields, the deployed policy succeeds on 97/175 scenarios, versus 65/175 before checking and 118/175 for the template (Table~\ref{tab:results}, panel A).
The policy--cascade difference is +18.3 percentage points (95\% interval 11.9 to 25.2); the policy--template difference is $-12.0$ points ($-17.9$ to $-6.4$).
The deployed policy retains 215/215 selected concept instances and has no detected polarity error, but still has 54 wrong-act outputs, four outputs with added content, and ten off-task outputs.
These diagnostics explain why complete retention alone is insufficient.

Supplying gold act and polarity raises the same deployed policy to 132/175, versus 71/175 before checking and 167/175 for the template (panel B).
The joint policy effect is +34.9 points (26.7 to 43.4); the gold-field--browser gap is +20.0 points (12.0 to 28.7).
The ungated gold-field cascade has 36 polarity errors, whereas the deployed output has zero detected errors.
Nevertheless, 37/108 accepted model outputs and 6/67 template fallbacks fail the full archived scorer.
Checker acceptance is therefore not interchangeable with independent success.

\subsection{Multilingual failure mechanisms}
The browser diagnostic sends 27 of 32 isolated feeling selections as requests.
Here 32 means eight feeling symbols in four locales, before any prompt-overlap exclusion; it is not the 30 feeling scenarios remaining in the primary score table.
The substring negation rule also mislabels three of 276 isolated symbol--language selections, including English \emph{Nose} and Gujarati \gu{નાક}.
Separating concept identity, act, and polarity directly addresses this input-layer defect.

The failure analysis includes all three Indic languages.
Recorded \lang{Hindi} feeling outputs include \bad{\hi{मुझे गुस्सा चाहिए।}} (\emph{“I want anger”}), where an intended state is realized as a request.
For \lang{Gujarati}, the eight apparent gold-field template failures are negated feelings with \gu{નથી}; the archived scorer omits this valid negator in its feeling-pattern test.
They are scorer lexical gaps, rather than evidence that these eight templates are wrong.
A recorded \lang{Marathi} tiredness output, \bad{\hi{मला दमलेला हवे आहे.}}, uses the request predicate \hi{हवे आहे} for an intended feeling; the archived scorer flags the act mismatch.
The checker must address act, perspective, and additions in each language; transferring English patterns alone does not do so.

The audit classifies the 37 accepted-model failures under gold fields as nine third-person reports, 12 negated-want constructions for feelings, ten malformed constructions, four additions, and two scorer gaps.
The six fallback failures are also scorer gaps.
These are assistant-produced diagnostic labels, not native-speaker judgments.
In particular, the lexical scorer's coverage affects both error counts and baseline rankings.

\subsection{Revised policy: regression evidence}
The repairs replace label-substring inference with concept-aware interpretation, add multilingual act and perspective checks, and check the fallback before returning it.
The same 175 scenarios informed these repairs and are therefore a regression set.
Under scorer v2, the archived policy scores 95/175 with browser fields and 138/175 with gold fields; the archived gold-field template scores 175/175 after its lexical gaps are corrected.
These are rescored archived sentences, not a new run of the old model.

The revised browser path and repaired template both score 175/175 on this set (Table~\ref{tab:results}, panel C).
The revised cascade alone scores 88/175, yielding a joint policy gain of 49.7 points (42.3 to 57.1).
This is a measured regression result, not evidence of generalization or superiority over the template.
A fresh 120-scenario set and blind rating materials exist, but native-speaker confirmation of their gold annotations and ratings are incomplete; we report no holdout performance or preference claim.
Figure~\ref{fig:workflow} and Table~\ref{tab:cases} show representative recorded cases rather than selecting only passing outputs.

\begin{table}[t]
\centering\footnotesize
\setlength{\tabcolsep}{2.8pt}
\begin{tabular}{@{}lrrrrrr@{}}
\toprule
&\multicolumn{3}{c}{English (reference)}&\multicolumn{3}{c}{Hindi}\\
\cmidrule(lr){2-4}\cmidrule(l){5-7}
Model & @1 & @5 & MRR & @1 & @5 & MRR \\
\midrule
Unigram & 3.6 & 11.6 & 0.069 & 6.4 & 14.5 & 0.099 \\
Bigram backoff & 13.6 & 29.5 & 0.204 & 24.0 & 43.1 & 0.320 \\
Trigram backoff & 16.0 & 32.8 & 0.234 & 25.9 & 44.0 & 0.334 \\
Suffix trie ($\infty$-gram) & 17.2 & 33.4 & 0.240 & 27.0 & 44.1 & 0.343 \\
Gutenberg HMM & 6.7 & 15.5 & 0.105 & -- & -- & -- \\
Hindi AAC HMM & -- & -- & -- & 24.5 & 41.9 & 0.315 \\
\midrule
\multicolumn{7}{@{}l}{Gujarati: no matched evaluation; shelf returns English.}\\
\multicolumn{7}{@{}l}{Marathi: no matched evaluation; reuses Hindi predictor.}\\
\bottomrule
\end{tabular}
\caption{Next-word prediction on AAC-user utterances from 60 of 300 source conversations per evaluated language (240 training; every fifth source held out). There are 2,220 English and 2,437 Hindi prediction contexts. Hit@1 and hit@5 are percentages; MRR is truncated at rank 10. Unigram is the frequency baseline. Dashes denote unmeasured cells, not zero performance.}
\label{tab:pred}
\end{table}


\subsection{Next-word prediction}
Table~\ref{tab:pred} evaluates AAC-user utterances in 60 held-out source conversations per evaluated language, not 60 utterances out of an unspecified pool.
Each language has 300 source conversations: 240 for training and 60 for testing, with every fifth source held out and no separate validation split in this evaluation.
All n-gram baselines and the evaluated suffix trie are rebuilt on the same 240 training sources, separately from the full-data deployed predictor.
Compared with trigram backoff, the trie gains 25/2,220 English and 29/2,437 Hindi hit@1 predictions: +1.126 points (95\% interval 0.505 to 1.766) and +1.190 points (0.463 to 1.922), computed from unrounded counts.
The corresponding hit@5 gains are 13/2,220 (+0.586 points) and 2/2,437 (+0.082 points).

Source conversations are disjoint, but 18 of 392 distinct English and 28 of 384 Hindi test utterances repeat training text.
Removing contexts from repeated utterances leaves 2,170 English and 2,352 Hindi contexts; trie--trigram hit@1 gains remain +0.968 and +1.020 points.
Out-of-vocabulary targets count as misses.
Gujarati and Marathi have no matched AAC prediction evaluation, so Table~\ref{tab:pred} records their coverage without inventing scores or treating language identity as a baseline method.

\subsection{Runtime and reproduction scope}
The revised browser-to-playback audit contains 32 runs across four languages, with one cold and 31 warm runs, on the GTX~1650.
Median click-to-first-playback is 9.9\,s and the 95th percentile is 20.8\,s; language medians are 12.4\,s (English), 15.3\,s (Hindi), 7.9\,s (Gujarati), and 10.7\,s (Marathi).
Every run plays a 440 Hz test tone: these are playback-integration timings, not click-to-audible-speech latency.
Greedy Indic retries repeated their first candidates in 40/40 uncached audit repeats, so the retry can add delay without changing the answer.

A CPU-only installation check reused cached weights; first-time download, clean GPU installation, and real synthesis were untested.
The model executes on the backend, but the browser requires externally hosted interface libraries and therefore is not verified as offline-ready.
No accessible public package or demo-video link was verified in the audit.
